\ssscp{\tsc{Sumba} phnology}{03-03-01-01}
Kambera has two low vowel segments with some overlap in their realisations.
They are transcribed by \citet{Klamer1998a} as ‹a› and ‹à›.
The vowel ‹a› is a tense vowel which is phonetically realised
as either short [a] or long [aː].
The vowel ‹à› is a lax vowel which is phonetically realised
as either [a] or [ɐ].
There is no clear rule or context which determines the realisations
(Marian Klamer pers. comm. March 2020).
We transcribe lax ‹à› as ‹ɐ›,
while we transcribe tense ‹a› as ‹a›.\footnote{
		The exact realisation of each allophone of Kambera low vowels
		is not specified in prose in \citet{Klamer1998a}.
		The allophones [a] and [aː] are probably slightly front,
		and [ɐ] is probably slightly back.
		\citet[23]{Klamer1998a} uses the symbols [a]
		and [ɑ] for the allophones of the lax vowel ‹à›.}
Based on examination of data in \citet{Onvlee1984},
it appears that some other languages of Sumba -- particularly
those of the \tsc{Central East Sumba} node (see \srf{03-03} for
details on the relationships of the languages of Sumba) -- also
have a contrast between two low vowels.
However, the lack of phonological descriptions means that we
cannot be sure of the phonetic realisations or how these vowels
fit into the phonological systems of the languages.

Languages of Sumba have vowel epenthesis (a.k.a. paragoge).
Certain words with a final consonant undergo epenthesis in order
to create a final open syllable.
The epenthetic vowel is extra-metrical
and does not shift stress from the penultimate vowel of the root.
In Kambera the epenthetic vowel is [u].
Examples include: *quzan > \it{uraŋ → uraŋ-u} ‘rain’,
*tasik > \it{tehik → tehik-u} ‘sea’,
and *umpu > \it{umbuk → umbuk-u} ‘grandchild’ \citep[19]{Klamer1998a}.
In Weyewa, Loli, and Mamboru the epenthetic vowel is usually
[a] (though exceptions occur); as in Weyewa \it{tazik-a} ‘sea’
and Mamboru \it{umbuk-a} ‘grandchild’.
In Kodi the epenthetic vowel is [o].

In some cases the final epenthetic vowel has become part of the
root and is only epenthetic from a historical perspective.
One example is Laboya \it{laŋta} ‘sky’ which is from PMP *laŋit
with final vowel epenthesis
and subsequent deletion of the preceding vowel.
The intermediate form with epenthesis
but no vowel deletion
is attested by Weyewa \it{laŋit-a} ‘sky’.\footnote{
		Similar patterns
		of final vowel epenthesis,
		often with subsequent deletion of the final vowel in the root,
		are also found in \tsc{Ambon-Seram} languages (\crf{ch:AS}).
		Examples include: *laŋit ‘sky’ > Haruku \it{lanit-o},
		Alune \it{lanit-e}, Yalahatan \it{lante},
		South Nuaulu \it{nante}, and Tamilou \it{laanto}.}

Some languages of western Sumba also have a contrast between
single and lengthened consonants.
The extent to which lengthening occurs
and the languages in which it occurs is currently unclear,
but sufficient data is available to confirm that it is a feature
of Kodi and Weyewa.
Other languages are likely to also have lengthened consonants.
Two minimal pairs demonstrating the contrast between single
and lengthened consonants in Kodi are /poto/ ‘four’ vs. /potːo/ ‘k.o.
bamboo’ and /loɗo/ ‘sing’ vs.
/loɗːo/ ‘day, sun’ \citep[36]{Lovestrand2021}.
Stressed vowels before single consonants in Kodi are phonetically
longer than those before lengthened consonants; /loɗo/ [ˈlɔˑɗɔ] ‘sing’ vs.
/loɗːo/ [ˈlɔʔɗɔ] ‘day, sun’.

In Laboya the difference between single
and lengthened consonants has been lost,
but the vowel length difference has become phonemic.
Examples include /taŋa/ ‘whole’ vs.
/taːŋa/ ‘enough’, /utu/ ‘fleas’ vs.
/uːtu/ ‘profit, advantage’
%(from Malay \it{untuŋ} ‘profit’),
and /moro/ ‘medicine’ vs. /moːro/ ‘green’ \citep[19]{Verdizade2019a}.
Orthographically, short vowels are often represented by analysts
and native speakers by doubling the following consonant; e.g.
‹uttu› /utu/ ‘fleas’ vs.
‹utu› \mbox{/uːtu/} ‘profit, advantage’
and ‹morro› /moro/ ‘medicine’ vs. ‹moro› /moːro/ ‘green’.

Information is conflicting as to whether Wanukaka has lengthened
consonants (like Kodi) or a vowel length difference (like Laboya),
and it is possible this language is undergoing a change from
one system to another.
Our main source of Wanukaka data,
\citet{Mitchell2008a}, has many instances of orthographic double
consonants, and \citet{MitchellND} describes this as representing “prolongation
of the consonant” or “a double or ‘geminate’ sound”.
In available recordings of Wanukaka \citep{LovestrandBalle2018},
some speakers do produce lengthened consonants
for words transcribed by \citet{Mitchell2008a} with a double
consonant, while other speakers do not.
The basis for this variation is currently unknown.

Similarly, our only source of Loli data,
\citet{Mitchell2008b}, has many instances of orthographic double consonants; e.g.
‹uttu› ‘louse’ vs. ‹utu› ‘luck’.
However, the lack of recordings or a phonological description
of Loli means that we cannot determine whether these double consonants
represent lengthened consonants or the length of the preceding vowel.
In this section, we treat orthographic double consonants in Wanukaka
and Loli as representing lengthened consonants.
It is also worth noting that \citet{Mitchell2008b} does not distinguish
between imploded and plain voiced stops in Loli.
If this representation is accurate,
Loli would be the only language of Sumba without two series of
voiced plosives (see \srf{03-03-02}).